\documentclass[10pt,a4paper]{amsart}
\usepackage[margin=19mm]{geometry}
\usepackage{amsmath,amssymb,mathtools}
\usepackage[scr=rsfs,cal=euler]{mathalfa}
\usepackage{xfrac}
\usepackage[unicode,hidelinks,bookmarks=false]{hyperref}
\newcommand{\ie}{i.e.\ }
\newcommand{\cf}{cf.\ }
\newcommand{\resp}{respectively\ }
\newcommand{\Subsection}{\subsection}
\providecommand{\nobreakdash}{\nobreak\hbox{-}}
\setlength{\emergencystretch}{2em}
\multlinegap0pt
\usepackage{amssymb}
\usepackage{mathtools}
\usepackage[shortlabels]{enumitem}
\usepackage{tikz}
\newtheorem{lemma}{Lemma}
\title{Independent Sets and Balanced Cycle-Linkings in 2-Connected Graphs}
\author{Shuaichao Wang}
\date{Source: arXiv:2609.15637v1 (CC BY 4.0)}
\begin{document}
\maketitle

\noindent\emph{Source and licence.} Independent Sets and Balanced Cycle-Linkings in 2-Connected Graphs. Version arXiv:2609.15637v1 (CC BY 4.0), distributed by arXiv under the Creative Commons Attribution 4.0 International licence, http://creativecommons.org/licenses/by/4.0/, as recorded in the arXiv licence metadata for this exact version and shown on the abstract page https://arxiv.org/abs/2609.15637v1. Full licence notice: \texttt{licenses/arxiv-2609.15637-CC-BY-4.0.txt}.\par\smallskip

\noindent\textit{Editorial scope.} Counted unit: the article's lemma lem:first-imbalance (source0.tex lines 723--735). The article's complete original proof of it is delivered with it (source0.tex lines 737--839, one \texttt{proof} environment of 103 lines). No mathematics is written, repaired or re-derived: the statement and the proof are the article's own lines, in the article's own notation, and no retained line is substituted or re-flowed.  No further source-local context is needed: every cross-reference of every retained line resolves inside the two delivered spans.  Nothing was omitted from any retained span, so the package carries no omission record.  No retained line contains a citation, so the wrapper carries no bibliography and no bibliography entry.  No retained line contains a figure environment, an image-inclusion command or drawing code, so no image is delivered and the package declares no asset dependency.  Editorial formatting only, in addition: an AMS \texttt{amsart} preamble that restates the article's own declarations and macros used by the retained lines, namely the environments \texttt{lemma} on separate counters, together with the standard packages those lines need.  The retained environments print in this document as Lemma 1 in order of appearance, whereas the article prints the same items with the numbering of its own full document.  The complete native source of the article as distributed in the arXiv e-print for this exact version is bundled unchanged as \texttt{sources/210435/source0.tex}.
\begin{lemma}\label{lem:first-imbalance}
Let $w$ be a nonbalanced cyclic binary word of length $\alpha$, and
let $r$ be the least length at which balance fails. Then
$
2\le r\le\lfloor\alpha/2\rfloor,
$
and $w$ has disjoint cyclic occurrences of $0z0$ and $1z1$, where
$z=z^R$ and $|z|=r-2$. Consequently, for some possibly empty words
$A,B$,
\begin{equation}\label{eq:obstruction-decomposition}
w\equiv 0z0\,A\,1z1\,B\equiv (z0A1z)(1B0).
\end{equation}
\end{lemma}

\begin{proof}
Length-one factors are balanced, so $r\ge2$. Every factor length can
be written as $a\alpha+s$, where $0\le s<\alpha$. Such a factor has
weight $a|w|_1$ plus the weight of its remaining length-$s$ factor.
Hence any imbalance at length $a\alpha+s$ already occurs at length
$s$, and therefore $r<\alpha$. If $r>\alpha/2$, the complementary
intervals of two length-$r$ factors witnessing an imbalance have
length $\alpha-r<r$ and the same absolute weight difference,
contradicting the minimality of $r$. Thus
$
2\le r\le\lfloor\alpha/2\rfloor.
$

Choose length-$r$ factors
$
u=u_1\cdots u_r,$$v=v_1\cdots v_r
$
such that
$
\delta:=|v|_1-|u|_1\ge2.
$
Since $r$ is the least imbalance length, the length-$(r-1)$ factors
obtained by deleting the first letters have weights differing by at
most one. Hence
\[
\delta-(v_1-u_1)
=
|v_2\cdots v_r|_1-|u_2\cdots u_r|_1
\le1.
\]
Since $\delta\ge2$, this gives $v_1-u_1\ge1$. As
$u_1,v_1\in\{0,1\}$, necessarily
$
u_1=0,v_1=1.
$
Substituting back gives $\delta\le2$, and hence $\delta=2$.
Applying the same argument after deleting the last letters yields
$
u_r=0,v_r=1.
$
Thus
$
u=0a0, v=1b1,
$
where $|a|=|b|=r-2$ and $|a|_1=|b|_1$.

We next show that $a=b$. Suppose otherwise, and let $p$ be their
maximal common prefix. If
$
a=p0a', b=p1b',
$
then $0p0$ and $1p1$ are equal-length factors of $w$ whose weights
differ by two, contradicting the minimality of $r$. If instead
$
a=p1a', b=p0b',
$
then $|a|_1=|b|_1$ implies
$
|b'|_1=|a'|_1+1.
$
Consequently, the equal-length factors $a'0$ and $b'1$ have weights
differing by two, again at a length smaller than $r$. Therefore
$a=b=:z$.

We claim that $z$ is a palindrome. Otherwise, comparing $z$ from
the two ends up to the first mismatch gives either
$
z=p0c1p^R
$ or $
z=p1c0p^R
$
for some words $p,c$. In the first case, the prefix $0p0$ of $0z0$
and the suffix $1p^R1$ of $1z1$ have weights differing by two; in
the second case, the same is true of $1p1$ and $0p^R0$. In either
case these factors have length $|p|+2<r$, contradicting the
minimality of $r$. Hence $z=z^R$.

It remains to show that the occurrences of $0z0$ and $1z1$ are
disjoint. They clearly cannot coincide. Suppose they overlap
properly. Since $2r\le\alpha$, there is a position outside their
union; cutting the cyclic word there turns the two occurrences into
overlapping linear intervals. Thus, for some $1\le h<r$, a suffix
of one occurrence equals a prefix of the other.

Suppose first that a length-$h$ suffix of $0z0$ equals a length-$h$
prefix of $1z1$. The former has weight equal to that of the last
$h-1$ letters of $z$, whereas the latter has weight one more than
that of the first $h-1$ letters of $z$. Since $z=z^R$, these two
parts of $z$ have the same weight, a contradiction. The opposite
overlap is excluded in the same way. Hence the two occurrences are
disjoint.

Their disjointness gives
$
w\equiv 0z0\,A\,1z1\,B
$
for some possibly empty words $A,B$. Rotating the cyclic word by
moving its initial $0$ to the end gives
$
w\equiv (z0A1z)(1B0),
$
which proves \eqref{eq:obstruction-decomposition}.
\end{proof}

\end{document}
